A MoLIC representa o único papel de usuário, quem faz o pedido \cite{molic}. As falas u: de uma cena são as ações disponíveis ao executor naquele estado; a fala d: seguinte é o conteúdo da próxima observação. A conclusão (d: OK) é condição de estado avaliada pelo verificador externo, nunca pela declaração do agente. Como a MoLIC não é modelo formal processável, o executor processa o statechart da Fig.~\ref{fig:estados}, no qual a ajuda é um subdiálogo que devolve a pessoa à cena de origem. A Tabela~\ref{tab:molic} dá as três leituras de cada elemento.
\begin{figure}[ht]
\centering
\setlength{\unitlength}{1pt}
\begin{picture}(238,112)
\put(0,87){\framebox(50,18){S0 · Início}}
\put(94,87){\framebox(58,18){S1 · Cardápio}}
\put(188,87){\framebox(48,18){S2 · Item}}
\put(50,98){\vector(1,0){44}}\put(59,101){\scriptsize abrir}
\put(94,91){\vector(-1,0){44}}\put(57,80){\scriptsize início}
\put(152,98){\vector(1,0){36}}\put(154,101){\scriptsize escolher}
\put(188,91){\vector(-1,0){36}}\put(157,80){\scriptsize incluir}
\put(94,40){\framebox(58,18){S3 · Revisão}}
\put(188,40){\framebox(48,18){S4 · Fim}}
\put(123,87){\vector(0,-1){29}}\put(127,69){\scriptsize rever}
\put(113,58){\vector(0,1){29}}\put(64,68){\scriptsize continuar}
\put(152,48){\vector(1,0){36}}\put(151,53){\scriptsize confirmar}
\put(212,87){\line(0,-1){24}}\put(212,63){\line(-1,0){73}}\put(139,63){\vector(0,-1){5}}\put(172,66){\scriptsize salvar}
\put(25,87){\line(0,-1){39}}\put(25,48){\vector(1,0){69}}\put(30,51){\scriptsize repetir ou rever}
\put(0,4){\framebox(101,18){Origem · S0 a S4}}
\put(158,4){\framebox(78,18){S5 · Ajuda}}
\put(101,17){\vector(1,0){57}}\put(110,21){\scriptsize consultar}
\put(158,9){\vector(-1,0){57}}\put(116,0){\scriptsize fechar}
\end{picture}
\caption{Transições principais da implementação. Editar uma linha em S3 reabre S2. Voltar preserva a origem da revisão ou da edição. Confirmar exige comanda não vazia. Uma fala inválida mantém a cena; o suporte restaura a origem.}
\label{fig:estados}
\end{figure}

\begin{table*}[!t]
\caption{Leituras de cada elemento da MoLIC}
\label{tab:molic}
\centering\scriptsize\setlength{\tabcolsep}{3pt}\renewcommand{\arraystretch}{1.1}
\begin{tabularx}{\textwidth}{@{}>{\raggedright\arraybackslash}p{.11\textwidth}>{\raggedright\arraybackslash}X>{\raggedright\arraybackslash}X>{\raggedright\arraybackslash}p{.25\textwidth}@{}}
\hline \textbf{Elemento MoLIC} & \textbf{Leitura da persona} & \textbf{Leitura do agente (observação ou ação)} & \textbf{Evento registrado}\\ \hline
Cena & Reconhece o assunto (início, cardápio, item, revisão, fim, ajuda) pelos signos da tela. & Observação: título e signos da cena que cruzam o recorte. & \UXev{tela_alterada} com cena anterior e nova\\
Fala u: & Solicita uma operação, como abrir um item, alterar a quantidade ou confirmar. & Ação: um gesto permitido sobre um alvo exposto, escolhido entre as alternativas do estado. & \UXev{acao_proposta}, \UXev{acao_executada}; proposta inválida: \UXev{acao_rejeitada}\\
Fala d: & Vê a resposta do preposto, como a nova quantidade da linha e o novo total. & Observação seguinte: o fragmento exposto da resposta. & \UXev{elemento_visto} e o evento de mutação correspondente\\
Processo do sistema & Percebe o efeito: comanda copiada, recalculada ou registrada. & Observação do efeito expresso; o estado integral fica com o verificador. & \UXev{item_adicionado}, \UXev{item_alterado}, \UXev{item_removido}, \UXev{pedido_confirmado}\\
Fala de recuperação & Volta aos detalhes, recompõe a linha, remove ou desfaz a remoção. & As mesmas ações de recuperação, quando expostas. & \UXev{item_alterado}, \UXev{item_removido}, \UXev{item_adicionado}\\
Acesso ubíquo à ajuda & Abre a ajuda em qualquer cena, escolhe a camada e fecha para voltar à origem. & Ação de protocolo sempre disponível; o conteúdo só chega após a abertura. & \UXev{ajuda_aberta} e \UXev{ajuda_fechada} com camada, tópico, origem e retorno\\ \hline
\end{tabularx}
\end{table*}

